%% LyX 2.4.4 created this file.  For more info, see https://www.lyx.org/.
%% Do not edit unless you really know what you are doing.
\documentclass[english]{article}
\usepackage[T1]{fontenc}
\usepackage[latin9]{inputenc}
\usepackage{enumitem}
\usepackage{amsmath}
\usepackage{amsthm}
\usepackage{cancel}
\usepackage{geometry}
\geometry{verbose,tmargin=1in,bmargin=1in,lmargin=1in,rmargin=1in}
\PassOptionsToPackage{normalem}{ulem}
\usepackage{ulem}

\makeatletter
%%%%%%%%%%%%%%%%%%%%%%%%%%%%%% Textclass specific LaTeX commands.
\numberwithin{equation}{section}
\numberwithin{figure}{section}
\newlength{\lyxlabelwidth}      % auxiliary length 

%%%%%%%%%%%%%%%%%%%%%%%%%%%%%% User specified LaTeX commands.
\usepackage{fancyhdr}
\usepackage{lastpage}
\pagestyle{fancy}
\usepackage{enumitem}
\usepackage{ifthen}
\everymath=\expandafter{\the\everymath\displaystyle}
%\DeclareMathSizes{10}{10}{10}{10}
\usepackage{multicol}

\usepackage{tikz}
\usetikzlibrary{patterns}
\usetikzlibrary{plotmarks}
\usepackage{pgfplots}
\pgfplotsset{compat=newest}

\usepackage{advdate}    % Advancing/saving dates
\usepackage{datetime}   % Dates formatting
\usepackage{datenumber} % Counters for dates
\usdate
% Added by lyx2lyx
\setlength{\parskip}{\smallskipamount}
\setlength{\parindent}{0pt}

\makeatother

\usepackage{babel}
\begin{document}
\renewcommand{\author}{Dr.\,James Ashe}

\newcommand{\classNum}{335}

\newcommand{\classSec}{A}

\newcommand{\examType}{Final Exam Review}

\renewcommand{\date}{\formatdate{15}{4}{2014}}

%%First page's header%%%%%%%%%%%%%%%%%%%%%%
\fancypagestyle{firststyle} %%header settings for first pagr
{
	\fancyhead[L]{MTH \classNum{}(\classSec{}) \author{}\\
	\textbf{\examType{}} \date{}}
	\fancyhead[C]{\textbf{}}
	\setlength{\headsep}{14pt}
}
%%%%%%%%%%%%%%%%%%%%%%%%%%%%%%%%%%%%%%%%%%

%%Next page's header%%%%%%%%%%%%%%%%%%%%%%%
\setlist{nolistsep}
\lhead{\classNum{}(\classSec{})} 
\chead{\textbf{}} 
\cfoot{\thepage{}~of~\pageref{LastPage}} 
\setlength{\headsep}{5pt} 
%%%%%%%%%%%%%%%%%%%%%%%%%%%%%%%%%%%%%%%%%%%

%%Various settings; see the preamble for other settings%%%%%
%\setenumerate[0]{leftmargin=12pt,itemindent=0pt,labelwidth=0pt}
\setlist{topsep=.5pt}
\renewcommand{\labelenumi}{\textbf{\arabic{enumi}.}}
\renewcommand{\labelenumii}{\textbf{\alph{enumii}.}}
\thispagestyle{firststyle}
%%%%%%%%%%%%%%%%%%%%%%%%%%%%%%%%%%%%%%%%%%

\newcommand{\dspace}{\vspace{-.15cm}}

\textbf{Please review the syllabus in regards to make-ups and tardiness
for exams. }Stuff you should know:\small

\subsubsection*{Sets\dspace}
\begin{itemize}
\item Basic notation/operations, e.g., $A\cap B$, $A\cup B$, $A-B$, $A\times B$
etc.

\begin{itemize}
\item Note: it makes no sense to apply these operations to elements. For
example, we can't write $a\cap A$ if $a$ is an element in $A$.
We \uline{can }write $\left\{ a\right\} \cap A$. 
\end{itemize}
\item To show that two sets are equal, $A=B$ show $A\subseteq B$ and $B\subseteq A$.

\begin{itemize}
\item To show that $A\subseteq B$, let $a\in A$ and then show that $a\in B$.
\end{itemize}
\item If $\left|A\right|=n$ then $A$ has $2^{n}$ subsets, including $\emptyset$.
\end{itemize}

\subsubsection*{Relations\dspace}
\begin{itemize}
\item Important definitions: reflexive, symmetric, transitive, and equivalence:
\item To show that $\mathcal{R}$ is one of these, apply the definition
to arbitrary $x,y,z$.
\item To show that $\mathcal{R}$ is not one of these (usually) use a counterexample.
\end{itemize}

\subsubsection*{Functions\dspace}
\begin{itemize}
\item Important definitions: one-to-one, onto, bijection, and group homomorphism.
\item Show a function is/isn't one-to-one, onto, or bijective.
\item Pigeon-hole Theorem.
\item Function composition. 
\end{itemize}

\subsubsection*{Groups\dspace}
\begin{itemize}
\item Important definitions: group, subgroup, and cyclic.
\item Prove/disprove that a set is a subgroup.
\item Apply Lagrange's Theorem.
\item You will not have to prove anything involving cosets. 
\item Review the HW results on the groups $S_{n}$ and $D_{n}$.
\item Find nontrivial subgroups of a given group.
\end{itemize}

\subsubsection*{Rings\dspace}
\begin{itemize}
\item Important definitions: rings
\item Show that a given set is a ring.

\begin{itemize}
\item If it is a known abelian subgroup, just state that and then show multiplicative
associativity and distribution.
\end{itemize}
\end{itemize}

\subsubsection*{Some general notes\dspace}
\begin{itemize}
\item Review HW problems. many of the exam questions are taken or modeled
after HW problems.
\item There will be some multiple choice problems. 

\begin{itemize}
\item [ex.]Let $G$ be an abelian group and $H$ be a subgroup of $G$.
Which of the following statements are correct?\\
\textbf{(a) }$H$ is cyclic.\\
\textbf{(b) $H$ }is an abelain group.\\
\textbf{(c) }$H$ is a finite group.\\
\textbf{(d) $H$ }is a non-abelian group.
\end{itemize}
\end{itemize}

\end{document}
